\section{Justificaci\'on}

La elecci\'on de Unity WebGL sobre librer\'ias nativas de JavaScript (Three.js, Babylon.js) o soluciones de \textit{streaming} remoto (Unreal Pixel Streaming) responde a requisitos concretos de visualizaci\'on t\'ecnica, desempe\~no web y viabilidad del \textit{pipeline} de arte t\'ecnico. El objetivo del proyecto no es solo mostrar un modelo 3D, sino construir un prototipo navegable, mantenible y evaluable acad\'emicamente, con interacciones t\'ecnicas, modos anal\'iticos y documentaci\'on consistente.

\subsection{Criterios t\'ecnicos de selecci\'on}

\subsubsection{Ejecuci\'on de l\'ogica en WebAssembly}

Unity utiliza IL2CPP para transpilar c\'odigo C\# a C++ y posteriormente a WebAssembly (WASM). Esto permite ejecutar la l\'ogica de interacci\'on y orquestaci\'on del prototipo en un entorno compilado y predecible dentro del navegador, con una arquitectura m\'as robusta para sistemas con m\'ultiples managers, estados de interfaz y herramientas de visualizaci\'on t\'ecnica.

\subsubsection{\textit{Pipeline} integrado de importaci\'on y renderizado}

A diferencia de un enfoque puramente manual en bibliotecas de bajo nivel, Unity ofrece un entorno integrado para importar geometr\'ia, gestionar materiales, perfilar rendimiento y desplegar un build WebGL desde una misma base de trabajo. Esta integraci\'on es especialmente valiosa cuando el proyecto requiere coordinar limpieza de activos CAD, shaders personalizados, UI responsiva y pruebas funcionales sin fragmentar el flujo entre demasiadas herramientas.

\subsubsection{Capacidad de extensi\'on para visualizaci\'on t\'ecnica}

El proyecto exige combinar renderizado realista con modos de inspecci\'on y an\'alisis. Unity URP permite aprovechar un flujo PBR consolidado y, al mismo tiempo, extender el sistema con shaders personalizados para corte transversal, vista \textit{X-Ray}, visualizaci\'on t\'ermica y otros modos anal\'iticos. Esta flexibilidad es determinante porque la versi\'on final del prototipo no depende de una sola visualizaci\'on, sino de una familia de modos con prop\'ositos distintos.

\subsubsection{Herramientas para colaboraci\'on entre programaci\'on y arte t\'ecnico}

El flujo de trabajo requiere colaboraci\'on entre programadores y artistas t\'ecnicos. Unity facilita esta articulaci\'on mediante materiales editables, inspecci\'on visual del contenido, herramientas de editor y un ecosistema orientado a iteraci\'on r\'apida. Esto reduce fricci\'on en tareas como ajuste de mallas, configuraci\'on de materiales, depuraci\'on de escenas y afinaci\'on de la experiencia visual.

\subsubsection{Escalabilidad y pertinencia disciplinar}

El uso de un motor ampliamente adoptado permite que el prototipo sea mantenible, replicable y escalable. Desde la perspectiva de Ingenier\'ia Multimedia, esta decisi\'on tambi\'en es pertinente porque conecta modelado 3D, optimizaci\'on gr\'afica, dise\~no de interfaz, evaluaci\'on con usuarios y publicaci\'on web en un mismo caso de estudio.

\subsubsection{Ventana de oportunidad tecnol\'ogica}

La convergencia entre WebAssembly, WebGL 2.0 y \textit{pipelines} modernos de renderizado abre una oportunidad concreta para democratizar el acceso a visualizaciones t\'ecnicas interactivas desde el navegador. El valor del proyecto no radica solo en optimizar recursos, sino en demostrar que una aplicaci\'on web puede soportar un caso de visualizaci\'on t\'ecnica con criterios reales de claridad espacial, usabilidad y rendimiento.

\subsection{Aporte a la Ingenier\'ia Multimedia}
\begin{itemize}
    \item documentar un \textit{pipeline} replicable para llevar activos CAD a un visor WebGL interactivo;
    \item validar el rol del ingeniero multimedia como puente entre complejidad t\'ecnica y experiencia de usuario;
    \item generar criterios verificables de rendimiento y consistencia funcional para visualizaci\'on web 3D;
    \item articular fundamentos de computaci\'on gr\'afica, UX y carga cognitiva en un proyecto aplicado.
\end{itemize}

\newpage
